\section{A relative Morse model in twelve dimensions}
\label{sec:relative-model}

We now turn a short generating list for an augmentation ideal into a Morse
function with few critical points.  Let $G$ be a finite group, put
$\Lambda=\Z[G]$, and let $I\subset\Lambda$ be its right augmentation ideal.
Assume that $I$ has a generating list of length $b\geq2$.
We will use this list and the elementary-row lemma established above.
Let $X$ be a closed connected oriented smooth four-manifold with
$\pi_1(X)\cong G$.  Since $G$ is finite, rational homology and Poincar\'e
duality give
\[
 H_1(X;\Q)=H_3(X;\Q)=0,
 \qquad \chi(X)=2+\dim_{\Q}H_2(X;\Q)\geq 2.
\]
No symplectic structure is needed for this part of the argument.
The relation between quadratic stabilization, group-ring complexes and
augmentation generators is developed by Eliashberg and Gromov
\cite[Sections~0.2.2 and~0.2.5]{EliashbergGromov1998}.
We carry out the required realization on a specified relative cobordism,
so that the number of variables and the exterior function remain fixed.

For $u,v\in\R^4$, put
\begin{equation}
 Q(u,v)=\frac{|v|^2-|u|^2}{2},
 \qquad r^2=|u|^2+|v|^2.
 \label{eq:quadratic}
\end{equation}
We also write $Q$ for its pullback to $X\times\R^8$.

\begin{proposition}
\label{prop:relative-morse}
Let $G$ be a finite group whose right augmentation ideal in $\Z[G]$
has a generating list of length $b\geq2$.
Let $X$ be a closed connected oriented smooth four-manifold with
$\pi_1(X)\cong G$.
There is a smooth Morse function $F:X\times\R^8\to\R$ such that
\[
 F=Q\quad\text{for }r\geq 3,
 \qquad
 \#\Crit(F)=\chi(X)+4b.
\]
Its critical points have indices between four and eight.
\end{proposition}

The number of added variables is part of the assertion.  We will prove the
proposition by changing handles on one fixed compact twelve-dimensional
cobordism.  This section constructs the cobordism and its handle complex;
Section~\ref{sec:cancellation} performs the reduction and completes the proof.

\subsection{The domain and its collars}

Define
\begin{equation}
 \begin{gathered}
 E=\{(u,v):r^2\leq 9,\ -1\leq Q\leq 1\},
 \qquad E_{\pm}=E\cap\{Q=\pm1\},\\
 W=X\times E,\qquad A=X\times E_-,\qquad B=X\times E_+.
 \end{gathered}
 \label{eq:relative-domain}
\end{equation}
The incoming and outgoing faces are $A$ and $B$; the side is the face
$r=3$.  At its corners, both $u$ and $v$ are nonzero and $dr,dQ$ are
linearly independent.  In a neighborhood of the side, the coordinates
\[
 (r,Q,u/|u|,v/|v|)
\]
give a product collar with interval coordinate $Q$.  Since $dQ\neq0$ on
$A\cup B$, these faces also have collars with coordinate $Q$.  To make
the collars compatible, choose a vector field with $dQ$-value one which
is the product coordinate field near the side, and extend it over the
ends by a partition of unity.  All handle constructions below take place
away from a fixed smaller side collar.

Radial contraction preserves $E$, so $E$ is contractible.  On $E_-$,
the parametrization
\[
 (\xi,v)\longmapsto
 \bigl(\sqrt{2+|v|^2}\,\xi,v\bigr),
 \qquad \xi\in S^3,\quad |v|^2\leq\frac72,
\]
identifies $E_-$ with $S^3\times D^4$.  Interchanging $u$ and $v$ gives
the same description of $E_+$.  Consequently both inclusions
$A\hookrightarrow W$ and $B\hookrightarrow W$ induce isomorphisms on
fundamental groups.

Write $\widetilde X$ for the universal cover of $X$ and use tildes for
the corresponding covers of $W,A,B$.  The relative homology of
$(E,E_-)$ is $\Z$ in degree four and zero in every other degree.  The
relative K\"unneth formula therefore gives, equivariantly for the deck
action,
\begin{equation}
 H_{4+i}(\widetilde W,\widetilde A;\Z)
 \cong H_i(\widetilde X;\Z).
 \label{eq:relative-homology}
\end{equation}
There is no torsion term because the relative homology of $(E,E_-)$ is
free.  The same formula holds with $\widetilde B$ in place of
$\widetilde A$.  In particular,
\begin{equation}
 H_4(\widetilde W,\widetilde A;\Z)\cong\Z,
 \qquad H_5(\widetilde W,\widetilde A;\Z)=0,
 \label{eq:low-relative-homology}
\end{equation}
where the action on $\Z$ is trivial.  Here we used connectedness and
simple connectedness of $\widetilde X$.

\subsection{An initial height function}

Choose a Morse function $f:X\to\R$ with one minimum and one maximum.
For completeness, this follows from the usual handle description of a
connected manifold.  After ordering the handles, the zero- and
one-handles form a connected graph.  Cancelling the edges of a spanning
tree against all but one zero-handle leaves one minimum.  Attachments
above each cancelling pair are transferred through the resulting
boundary identification.  Applying the same procedure to the dual
handle decomposition removes the extra maxima without changing the
minimum.  This uses the handle decomposition theorem and the geometric
cancellation criterion, in their elementary zero/one-handle cases
\cite[Theorem~2.22 and Lemma~8.24]{Ranicki2002}.

Let $\eta:[0,\infty)\to[0,1]$ be smooth, equal to one near zero and
supported in $[0,1)$.  For $\sigma>0$ sufficiently small, set
\begin{equation}
 h_0(x,u,v)=Q(u,v)+\sigma\eta(r^2)f(x).
 \label{eq:initial-height}
\end{equation}
Choose $\sigma$ so that
\[
 \sigma\|f\|_{\infty}<\frac12,
 \qquad
 2\sigma\|f\|_{\infty}\|\eta'\|_{\infty}<1.
\]
The perturbation is supported where $r<1$ and hence $|Q|<1/2$.
Thus $h_0(W)\subset[-1,1]$, and $h_0=Q$ near every face of $W$.
For $r>0$, its derivative in the extra variables has norm at least
\[
 r\bigl(1-2\sigma\|f\|_{\infty}\|\eta'\|_{\infty}\bigr)>0.
\]
At $r=0$, the critical-point equation is $df=0$, and the Hessian is
the direct sum of $\sigma\operatorname{Hess}f$ and
$\operatorname{Hess}Q$.  Therefore $h_0$ is Morse, with exactly the
critical points of $f$, each shifted upward by four in index.  In
particular it has one critical point of index four and one of index
eight.

\subsection{Levels and the relative handle complex}

Use gradient-like fields which are product fields on the side collar.
The Morse--handle correspondence then describes $(W;A,B)$ by handles
attached in the interiors of its regular levels.  The standard local
handle traces and their product regions give this correspondence
relative to the side collar; see
\cite[Proposition~2.20 and Theorem~2.22]{Ranicki2002} for the handle
construction.  Only handles of indices $4,\ldots,8$ will occur.

We may order these handles by index.  Indeed, if an index-$j$ handle
is followed by an index-$\ell$ handle with $\ell\leq j$, the latter
attaching sphere has dimension $\ell-1$, and the former belt sphere
has dimension $11-j$ in the intervening eleven-dimensional level.
Their dimension sum is
\[
 (\ell-1)+(11-j)=10+\ell-j<11.
\]
General position makes the attaching sphere disjoint from that belt
sphere.  After shrinking its framed neighborhood, it can be moved
back through the surgery, interchanging the two handles.  All later
attachments are transferred through the induced identification of
the upper level.  The same argument permits arbitrary ordering
within an equal-index block, with all moves supported in the interior.

Every regular level $L$ is connected, and its inclusion into $W$
induces
\begin{equation}
 \pi_1(L)\cong\pi_1(W)\cong G.
 \label{eq:level-fundamental-group}
\end{equation}
To see this directly, surgery for an index-$j$ handle removes a
neighborhood of $S^{j-1}$ and inserts one of $S^{11-j}$.  For
$4\leq j\leq8$, both spheres have dimension at least three and
codimension at least four in the level.  General position and van
Kampen show that these surgeries preserve connectedness and the
fundamental group.  Start with $A$, whose inclusion into $W$ is a
fundamental-group isomorphism.  Alternatively, from any intermediate
level both the upper part and the reversed lower part are built with
handles of index at least four, so neither part changes the
fundamental group.  The interior of $L$ has the same fundamental
group as $L$.

Choose an orientation and a lift of each handle.  Regard the deck
action as a right action: if it was originally written on the left
as $D_g$, use $c\cdot g=D_{g^{-1}}c$.  Let $C_i$ be the based free
right $\Lambda$-module generated by the handles of index $4+i$, and
let $r_i$ be their number.  The resulting relative handle complex is
\begin{equation}
 0\longrightarrow C_4\xrightarrow{\partial_4}C_3
 \xrightarrow{\partial_3}C_2\xrightarrow{\partial_2}C_1
 \xrightarrow{\partial_1}C_0\longrightarrow0,
 \qquad C_i\cong\Lambda^{r_i}.
 \label{eq:handle-complex}
\end{equation}
Its degree-$i$ homology is
$H_{4+i}(\widetilde W,\widetilde A;\Z)$.  One way to obtain the
complex is to filter the handle decomposition by index; each
successive relative group is freely generated by the lifted handles
in that index.  Boundary coefficients are signed intersections of
lifted attaching spheres with lifted belt spheres one index lower.
Thus, writing
\[
 \partial e_j=\sum_i e_i a_{ij},
\]
the entries $a_{ij}$ lie in $\Lambda$ and multiply the basis elements
on the right.  Replacing a basis vector $e_j$ by $e_j+e_i q$, with $i\neq j$, adds
column$_i\,q$ to column$_j$ in its outgoing differential.  If $T$
is the elementary matrix of this basis change, the adjacent
differentials change coherently as $D\mapsto DT$ and
$D'\mapsto T^{-1}D'$.  These conventions will be used for the handle
slides in Section~\ref{sec:cancellation}.

\begin{lemma}
\label{lem:handle-complex}
For any relative handle presentation of this same $(W;A,B)$ with
indices in $4,\ldots,8$ and $r_0=r_4=1$, its based right-module
complex satisfies
\begin{equation}
 \operatorname{im}\partial_1=I,
 \qquad
 \operatorname{im}\partial_2=\ker\partial_1,
 \qquad
 r_2=\chi(X)-2+r_1+r_3\geq r_1,
 \label{eq:handle-identities}
\end{equation}
where $C_0$ is identified with $\Lambda$ using its chosen basis.
The same assertions hold for the reversed cobordism relative to $B$.
\end{lemma}

\begin{proof}
By \eqref{eq:low-relative-homology}, the quotient
$C_0/\operatorname{im}\partial_1$ is the trivial right module $\Z$.
A surjective right-module map $\Lambda\to\Z$ sends $1$ to $1$ or
$-1$, so its kernel is $I$.  The vanishing of relative homology in
degree five gives the second equality.

Since $E$ is contractible and $E_-$ has the homotopy type of $S^3$,
we have $\chi(W,A)=\chi(X)$.  All handle indices have been shifted
by the even number four, and hence
\[
 \chi(X)=r_0-r_1+r_2-r_3+r_4
         =2-r_1+r_2-r_3.
\]
This gives the rank identity; its inequality follows from
$\chi(X)\geq2$ and $r_3\geq0$.  Finally, the same relative covering
homology computation holds with $B$ in place of $A$.  Reversal sends
index $4+i$ to index $8-i$ and preserves the endpoint counts, so the
argument applies unchanged.
\end{proof}

We have now constructed a height on the required fixed cobordism
and identified the algebra that its handles must satisfy.  To finish
Proposition~\ref{prop:relative-morse}, it remains to reduce the
index-five and index-seven counts to $b$ while retaining the single
index-four and index-eight handles and the prescribed collars.
Then Lemma~\ref{lem:handle-complex} will force
$r_2=\chi(X)-2+2b$ and the total count $\chi(X)+4b$.
